\par\Needspace{4\baselineskip}
\begin{center}
\SourceLeaf{055}\pdfbookmark[2]{37.}{article-037}37.
\end{center}
Let us return to the equation relating the eccentric anomaly to the time in elliptic motion
\[
E-e\sin E=\frac{kt\BareRoot{(1+\mu)}}{a^{\frac{3}{2}}},
\]
where $E$ is assumed to be expressed in parts of the radius. From here on we shall omit the factor $\BareRoot{(1+\mu)}$. If a case should ever arise in which it is possible and worth the effort to take this factor into account, the symbol $t$ would denote the time elapsed since perihelion multiplied by $\BareRoot{(1+\mu)}$, in place of the elapsed time itself. Further, we denote the distance at perihelion by $q$ and introduce, in place of $E$ and $\sin E$, the quantities $E-\sin E$ and $E-\tfrac{1}{10}(E-\sin E)=\tfrac{9}{10}E+\tfrac{1}{10}\sin E$. The attentive reader will see from what follows why we choose these particular quantities. Our equation then takes the following form
\[
(1-e)(\tfrac{9}{10}E+\tfrac{1}{10}\sin E)+(\tfrac{1}{10}+\tfrac{9}{10}e)(E-\sin E)=kt\left(\frac{1-e}{q}\right)^{\frac{3}{2}}.
\]
If $E$ is regarded as a small quantity of the first order, then $\tfrac{9}{10}E+\tfrac{1}{10}\sin E=E-\tfrac{1}{60}E^3+\tfrac{1}{1200}E^5-\text{etc.}$ is a quantity of the first order, whereas $E-\sin E=\tfrac{1}{6} E^3-\tfrac{1}{120}E^5+\tfrac{1}{5040}E^7-\text{etc.}$ is a quantity of the third order. Thus, putting
\[
\frac{6(E-\sin E)}{\frac{9}{10}E+\frac{1}{10}\sin E}=4A,\qquad
\frac{\frac{9}{10}E+\frac{1}{10}\sin E}{2\BareRoot{A}}=B,
\]
we find that $4A=E^2-\tfrac{1}{30}E^4-\tfrac{1}{5040}E^6-\text{etc.}$ is a quantity of the second order and that $B=1+\tfrac{3}{2800}E^4-\text{etc.}$ differs from unity by a quantity of the fourth order. Our equation then becomes
\[
\text{1)}\qquad B\left\{2(1-e)A^{\frac{1}{2}}+\tfrac{2}{15}(1+9e)A^{\frac{3}{2}}\right\}=kt\left(\frac{1-e}{q}\right)^{\frac{3}{2}}.
\]

Ordinary trigonometric tables suffice to calculate $\tfrac{9}{10}E+\tfrac{1}{10}\sin E$ with adequate precision. They give inadequate precision for $E-\sin E$ whenever $E$ is a small angle, so this method would give insufficiently accurate values of $A$ and $B$. A special table would overcome this difficulty by allowing us to use $E$ as the argument and look up either $B$ itself or the logarithm of $B$. Anyone with even a modest knowledge of analysis can readily see how to construct such a table. Using the equation
\[
\frac{9E+\sin E}{20B}=\BareRoot{A}
\]
\SourceLeaf{056}we could also determine $\BareRoot{A}$, and hence $t$ by formula 1), with any desired precision.

Here is a specimen of such a table, which at least shows how slowly $\log B$ increases. It would be superfluous to extend the table, since below we shall describe tables in a much more convenient form.
\begin{center}
\begin{tabular}{|r|r||r|r||r|r|}
\hline
$E$ & $\log B$ & $E$ & $\log B$ & $E$ & $\log B$\\
\hline
$0^0$ & $0{,}000\,0000$ & $25^0$ & $0{,}000\,0168$ & $50^0$ & $0{,}000\,2675$\\
$5$ & $00$ & $30$ & $0349$ & $55$ & $3910$\\
$10$ & $04$ & $35$ & $0645$ & $60$ & $5526$\\
$15$ & $22$ & $40$ & $1099$ & & \\
$20$ & $69$ & $45$ & $1758$ & & \\
\hline
\end{tabular}
\end{center}
